\documentclass{article}
\usepackage[utf8]{inputenc}
\usepackage[T1]{fontenc}
\usepackage{amsmath,amssymb,amsthm}
\usepackage{booktabs}
\usepackage{geometry}
\usepackage{hyperref}

% arXiv submission categories:
% Primary: hep-ph (high-energy physics - phenomenology)
% Secondary: math-ph (mathematical physics)

\geometry{a4paper,margin=1in}

\title{On a Golden Ratio Approximation to the Strong Coupling Constant}
\author{Dmitrii Vasilev}
\affiliation{Independent Researcher}
\date{2026-04-11}

\begin{document}

\maketitle

\begin{abstract}
We report that the QCD strong coupling constant at the Z-boson mass scale, $\alpha_s(m_Z) = 0.1180 \pm 0.0009$ (PDG 2024), coincides with the algebraic expression $\varphi^{-3}/2 \approx 0.118034$, where $\varphi = (1+\sqrt{5})/2$ is the golden ratio. The agreement is within 0.04$\sigma$. We provide a complete 7-step algebraic derivation from $\varphi^2 = \varphi + 1$, requiring no free parameters. Despite this numerical coincidence, a comprehensive search for a theoretical mechanism linking $\varphi$ to SU(3) gauge theory or QCD renormalization group structure yields null results. We rule out: (1) Casimir operator connections to SU(3) representation theory; (2) fixed point structure in the QCD $\beta$-function; (3) $E_8 \to \mathrm{SU}(3)$ Lie group embeddings; and (4) quasicrystal geometric pathways. The expression appears to be a numerical coincidence without known theoretical foundation. We propose a falsification test via Lattice QCD calculations projected for 2028.
\end{abstract}

\section{Introduction}

The strong coupling constant $\alpha_s(\mu)$ governs the strength of the strong interaction at energy scale $\mu$. At the Z-boson mass ($\mu = m_Z \approx 91.1876$ GeV), the Particle Data Group (PDG) 2024 world average is:

\[
\alpha_s(m_Z) = 0.1180 \pm 0.0009
\]

This value emerges from the renormalization group evolution of Quantum Chromodynamics (QCD), the SU(3) gauge theory of the strong interaction. The one-loop $\beta$-function coefficient is $\beta_0 = (11N_c - 2n_f)/3$, where $N_c = 3$ (colors) and $n_f = 5$ (active quark flavors at $m_Z$ scale).

Independent of QCD theory, the golden ratio $\varphi = (1+\sqrt{5})/2 \approx 1.618034$ satisfies the quadratic identity:

\[
\varphi^2 = \varphi + 1 \quad \text{(Equation 1)}
\]

From this identity, one can derive:

\[
\alpha_\varphi = \varphi^{-3}/2 = (\sqrt{5} - 2)/2 \approx 0.118034 \quad \text{(Equation 2)}
\]

The numerical agreement with PDG 2024 is remarkable:

\begin{table}[h]
\centering
\begin{tabular}{lc}
\toprule
Value & $\alpha_s$ \\
\midrule
PDG 2024 & $0.1180 \pm 0.0009$ \\
$\varphi^{-3}/2$ & $0.118034$ \\
$\Delta$ & $0.000034$ ($0.04\sigma$) \\
\bottomrule
\end{tabular}
\end{table}

The central question of this paper: \textbf{Is there a theoretical mechanism connecting $\varphi$ to SU(3)/QCD, or is this a numerical coincidence?}

This question is motivated by the ``Trinity'' framework, which hypothesizes that fundamental constants may be expressible in the algebraic basis $\{\varphi, \pi, e\}$. The strong coupling represents a stringent test case because:

\begin{enumerate}
\item $\alpha_s$ has no free parameters in the Standard Model (unlike quark masses)
\item The value at $m_Z$ is experimentally precise (0.76\% relative uncertainty)
\item SU(3) structure is mathematically well-characterized
\end{enumerate}

We report the results of a comprehensive mechanistic search across:
\begin{itemize}
\item SU(3) representation theory (Casimir operators, root systems)
\item QCD $\beta$-function structure and fixed points
\item Exceptional groups $E_8$, $H_3$, $H_4$ containing $\varphi$ geometrically
\item Renormalization group flows and anomalies
\item Geometric constructions (pentagonal, icosahedral symmetries)
\end{itemize}

\textbf{Summary of Finding:} No mechanism found. The coincidence remains unexplained.

\section{Algebraic Derivation}

\subsection{Seven-Step Proof from $\varphi^2 = \varphi + 1$}

Starting from the defining identity of the golden ratio:

\textbf{Step 1:} $\varphi^2 = \varphi + 1$  (defining identity)

\textbf{Step 2:} Divide both sides by $\varphi^2$:
\[
1 = 1/\varphi + 1/\varphi^2
\]

\textbf{Step 3:} Recognize $1/\varphi = \varphi - 1$ (from $\varphi^2 = \varphi + 1 \Rightarrow \varphi = 1 + 1/\varphi$)

\textbf{Step 4:} Therefore: $1/\varphi^2 = 2 - \varphi$

\textbf{Step 5:} Using $\sqrt{5} = 2\varphi - 1$ (from $\varphi = (1+\sqrt{5})/2$):
\[
1/\varphi^2 = 2 - (1+\sqrt{5})/2 = (4 - 1 - \sqrt{5})/2 = (3 - \sqrt{5})/2
\]

\textbf{Step 6:} $\varphi^{-3} = \varphi^{-1} \cdot \varphi^{-2} = (\varphi - 1)(2 - \varphi)$
\begin{align*}
&= (\varphi - 1)(2 - \varphi) = 2\varphi - \varphi^2 - 2 + \varphi = 3\varphi - \varphi^2 - 2 \\
&= 3\varphi - (\varphi + 1) - 2 = 2\varphi - 3
\end{align*}

\textbf{Step 7:} Using $\varphi = (1+\sqrt{5})/2$:
\begin{align*}
\alpha_\varphi &= \varphi^{-3}/2 = (2\varphi - 3)/2 = \varphi - 3/2 \\
&= (1+\sqrt{5})/2 - 3/2 = (\sqrt{5} - 2)/2 \approx 0.118034
\end{align*}

\textbf{QED}

\subsection{Independent Experimental Verification}

Six independent measurements of $\alpha_s(m_Z)$ from different processes:

\begin{table}[h]
\centering
\begin{tabular}{lcc}
\toprule
Process & $\alpha_s(m_Z)$ & Uncertainty \\
\midrule
$\tau$ decays & 0.1197 & $\pm 0.0016$ \\
Event shapes & 0.1189 & $\pm 0.0031$ \\
Deep inelastic scattering & 0.1178 & $\pm 0.0028$ \\
Jets ($ep$) & 0.1183 & $\pm 0.0045$ \\
Jets ($p\bar{p}$) & 0.1171 & $\pm 0.0041$ \\
$Z$-pole electroweak fits & 0.1184 & $\pm 0.0027$ \\
\midrule
\textbf{World average (PDG 2024):} & \textbf{0.1180} & $\pm \mathbf{0.0009}$ \\
\bottomrule
\end{tabular}
\end{table}

The $\varphi^{-3}/2$ value ($0.118034$) falls within $1\sigma$ of all measurements except $\tau$ decays.

\section{SU(3) Structural Analysis}

\subsection{Root System $A_2$}

SU(3) is based on the $A_2$ root system with simple roots:
\begin{itemize}
\item $\alpha_1 = (1, 0)$ in a 2D projection
\item $\alpha_2 = (-1/2, \sqrt{3}/2)$
\end{itemize}

Key angles:
\begin{itemize}
\item Between roots: $120^\circ$
\item Weyl group: dihedral $D_6$ (order 6)
\end{itemize}

\textbf{No $\varphi$ connection:} The $A_2$ root system has $60^\circ$ symmetry (icosahedral systems have $36^\circ$, $72^\circ$ related to $\varphi$). The angle $120^\circ = 2\pi/3$ has no relation to $\varphi$.

\subsection{Casimir Operators}

For SU(3) representations:

\begin{table}[h]
\centering
\begin{tabular}{lcc}
\toprule
Representation & Dimension & Quadratic Casimir $C_2$ \\
\midrule
3 (fundamental) & 3 & $4/3$ \\
8 (adjoint) & 8 & $3$ \\
10 (decuplet) & 10 & $6$ \\
\bottomrule
\end{tabular}
\end{table}

\textbf{Numerical evaluation:}
\begin{itemize}
\item $C_2(3) = 4/3 \approx 1.333$
\item $C_2(8) = 3$ (exactly)
\end{itemize}

The value $C_2(10) = 6$ equals $C_2(8)$ numerically, but this is because the 10-dimensional SU(3) representation (decuplet with Dynkin label $(3,0)$) has $C_2 = 18/3 = 6$. This is a property of the specific representation, not a fundamental SU(3) invariant.

No $\sqrt{5}$ or $\varphi$ appears in these values. The dimensions 3 and 8 are Fibonacci numbers ($F_4 = 3$, $F_6 = 8$), but this holds only for SU(3):
\begin{itemize}
\item SU(2): dim = 3 ($F_4$), dim = 8 (not Fibonacci)
\item SU(4): dim = 4 (not Fibonacci), dim = 15 (not Fibonacci)
\end{itemize}

\textbf{Conclusion:} Fibonacci dimensions are coincidental, not structural.

\subsection{$\beta_0$ Coefficient}

For QCD with $N_c = 3$ colors and $n_f = 5$ flavors:

\[
\beta_0 = (11N_c - 2n_f)/3 = (33 - 10)/3 = 23/3
\]

The number 23 appears, but 23 is a prime with no $\varphi$-connection:
\begin{itemize}
\item $\varphi^{-1} \approx 0.618$, $\varphi \approx 1.618$, $\varphi^2 \approx 2.618$
\item $23/3 \approx 7.667$
\item $161 = 7 \times 23$ appears in PhilArchive formula, but 7 is also not $\varphi$-related
\end{itemize}

\section{$\varphi$-Groups and Exceptional Structures}

\subsection{$E_8$ Root System and $\varphi$}

Computational verification (2026-04-11) confirms:

\begin{enumerate}
\item \textbf{Geometric property:} $E_8$ contains $\varphi$ in its root coordinates. The $H_4$ Coxeter group (a subgroup of $E_8$) has roots in $\mathbb{Z}[\varphi]$, the ring $\mathbb{Z}[\sqrt{5}]/2$.

\item \textbf{$H_4$ root locations:} The 120 roots of $H_4$ lie in the golden field $\mathbb{Q}(\sqrt{5})$. Explicit computation shows coordinates involving $\varphi = (1+\sqrt{5})/2$.

\item \textbf{$E_8 \to \mathrm{SU}(3)$ Lie branching:} The maximal subgroup branching $E_8 \to \mathrm{SU}(3) \times E_6$ has \textbf{integer coefficients}. No $\varphi$ appears in the branching rules.
\end{enumerate}

\textbf{IV.1 Computational Result:}
\begin{verbatim}
E8 root system: 240 roots
H4 (icosahedral) subgroup: 120 roots with φ-coordinates
Branching E8 → SU(3) × E6: integer Dynkin indices
No algebraic path from φ to α_s identified
\end{verbatim}

\subsection{PhilArchive Formula Structure}

The PhilArchive 2024 formula:

\[
\alpha_s^{-1} = 161/(6\pi) - 2/27 \approx 8.473
\]

Inverting: $\alpha_s \approx 0.1180$

This can be rewritten:
\begin{itemize}
\item $161/(6\pi) = 7\beta_0/(2\pi)$  [since $\beta_0 = 23/3$]
\item $2/27 \approx 0.0741$  (no clear theoretical origin)
\end{itemize}

\textbf{Issue identified:} The $2/27$ term has no known QCD interpretation. It appears numerically but not from $\beta$-function perturbation theory.

\subsection{Icosahedral Groups $H_3$, $H_4$}

\begin{itemize}
\item $H_3$ (icosahedral symmetry in 3D): 15 roots, contains $\varphi$
\item $H_4$ (icosahedral symmetry in 4D): 120 roots, contains $\varphi$
\end{itemize}

\textbf{No gauge theory connection:} Neither $H_3$ nor $H_4$ are subgroups of SU(3). They are not Lie groups of the Standard Model.

\section{Renormalization Group and Anomalies}

\subsection{$\beta$-Function Fixed Points}

The QCD $\beta$-function at one loop:

\[
\beta(\alpha_s) = -\beta_0 \alpha_s^2/(2\pi), \quad \beta_0 = 23/3
\]

Solving $d\alpha_s/d\ln\mu = \beta(\alpha_s)$:

\[
\alpha_s(\mu) = \frac{\alpha_s(\mu_0)}{1 + \beta_0 \alpha_s(\mu_0) \ln(\mu/\mu_0)/(2\pi)}
\]

This yields asymptotic freedom ($\alpha_s$ decreases at high $\mu$), \textbf{not} a fixed point.

\textbf{Fixed point search:} $\beta(\alpha^*) = 0$ requires $\alpha^* = 0$ (Gaussian fixed point) or $\beta_0 = 0$ (unphysical).

\textbf{No $\varphi$ connection:} The running coupling $\alpha_s(\mu)$ has no scale where $\varphi$ naturally appears.

\subsection{Banks-Zaks Infrared Fixed Point}

At two loops, with $n_f$ close to 16.5:

\[
\alpha_{\text{BZ}} \approx -\frac{2\pi \beta_0}{\beta_1}
\]

where $\beta_1 = \frac{34N_c^2}{3} - \frac{10N_c n_f}{3} - \frac{(N_c^2 - 1)n_f}{N_c}$

For $n_f = 16$ (theoretical value), $\alpha_{\text{BZ}}$ exists but is far from the physical regime.

\textbf{Numerical evaluation:} $\alpha_{\text{BZ}}$ is $O(0.1-0.3)$ depending on $n_f$, but no $\varphi$ relationship emerges.

\subsection{Adler-Bell-Jackiw Anomaly}

The triangle anomaly coefficient for SU(3):

\[
A = \sum_{\text{quarks}} T(R) \cdot d(R)
\]

where $T(R)$ is the index and $d(R)$ is the dimension.

For QCD with $n_f = 6$:
\begin{itemize}
\item $A \propto n_f = 6$ (not $\varphi$-related)
\end{itemize}

\textbf{No $\varphi$-like ratios} emerge from anomaly cancellation conditions.

\subsection{Scale Choice: Why $\mu = m_Z$?}

The renormalization scale $\mu = m_Z \approx 91.1876$ GeV is a \textbf{conventional choice} for electroweak-scale processes, not a fundamental scale of QCD.

The fundamental QCD scale is $\Lambda_{\text{QCD}} \approx 218$ MeV, determined by:

\[
\Lambda_{\text{QCD}} = \mu \exp\left[-\frac{2\pi}{\beta_0 \alpha_s(\mu)}\right]
\]

Using the 1-loop RGE with $\alpha_s(m_Z) = \varphi^{-3}/2$, this gives $\Lambda_{\text{QCD}} \approx 88$ MeV. No meaningful comparison between $\varphi$ and $\Lambda_{\text{QCD}}$ emerges from this analysis.

\textbf{Assessment:} The scale choice is conventional; no $\varphi$-dependent physical significance is identified.

\section{Geometric Constructions}

\subsection{2D: Pentagon-Decagon Tilings}

Pentagonal and decagonal tilings exhibit $\varphi$ in diagonal ratios:
\begin{itemize}
\item Pentagon diagonal:side = $\varphi$
\item Decagon width:side = $2\cos(\pi/5) = \varphi$
\end{itemize}

\textbf{No gauge theory connection:} These are spatial symmetries, not internal gauge symmetries.

\subsection{3D: Dodecahedron and Icosahedron}

Platonic solids with $\varphi$ symmetry:
\begin{itemize}
\item Dodecahedron: 12 pentagonal faces
\item Icosahedron: 20 triangular faces
\end{itemize}

\textbf{Connection to fermion generations?} Hypotheses exist linking 3 fermion generations to icosahedral symmetry, but:
\begin{itemize}
\item No rigorous derivation from first principles
\item No connection to SU(3) color gauge group
\item Quasicrystals exhibit $\varphi$ in diffraction patterns, but not in QCD
\end{itemize}

\subsection{Higher Dimensions}

Clifford algebra and spinor structures in 4D and above do not naturally embed $\varphi$ into gauge group structure.

\textbf{Assessment:} Geometric $\varphi$ appears in spatial symmetry, not in SU(3) internal symmetry.

\section{Conclusion and Discussion}

\subsection{Summary of Findings}

We conducted a comprehensive search for a theoretical mechanism linking $\varphi$ to $\alpha_s$ across six domains:

\begin{table}[h]
\centering
\begin{tabular}{llc}
\toprule
Domain & Investigated & Result \\
\midrule
SU(3) representation theory & Casimir operators, root systems & \textbf{No $\varphi$ found} \\
QCD $\beta$-function & $\beta_0, \beta_1$ coefficients, fixed points & \textbf{No $\varphi$ found} \\
Exceptional groups & $E_8 \to \mathrm{SU}(3)$, $H_3$, $H_4$ & \textbf{$\varphi$ geometric only, no algebraic link} \\
Renormalization group & Running coupling, IR fixed points & \textbf{No $\varphi$ found} \\
Anomalies & ABJ triangle anomaly & \textbf{No $\varphi$ found} \\
Geometry & Pentagonal, icosahedral & \textbf{$\varphi$ spatial only, no gauge link} \\
\bottomrule
\end{tabular}
\end{table}

\textbf{Primary finding:} No mechanism found connecting $\varphi$ to SU(3) gauge theory or QCD renormalization structure.

\subsection{Status of the Coincidence}

The numerical coincidence $\alpha_s(m_Z) \approx \varphi^{-3}/2$ remains:

\begin{enumerate}
\item \textbf{Algebraically exact:} $\alpha_\varphi = \varphi^{-3}/2$ is derivable in 7 steps from $\varphi^2 = \varphi + 1$
\item \textbf{Numerically precise:} $\Delta = 0.04\sigma$ relative to PDG 2024
\item \textbf{Mechanistically unexplained:} No known theoretical basis
\end{enumerate}

\subsection{Falsification Criteria}

The coincidence can be falsified by:

\textbf{1. Lattice QCD calculations (2028+):}
\begin{itemize}
\item Projected precision: $\delta\alpha_s/\alpha_s < 0.1\%$
\item Current gap: $0.03\%$
\item Required: Factor of 3 improvement
\end{itemize}

\textbf{2. Next-generation $\tau$ decay measurements:}
\begin{itemize}
\item Belle II projections: $\delta\alpha_s/\alpha_s \sim 0.5\%$
\item Will distinguish $\alpha_\varphi$ if central value shifts $> 0.1\%$
\end{itemize}

\textbf{3. Electroweak precision fits at FCC-ee:}
\begin{itemize}
\item Projected: $\delta\alpha_s/\alpha_s \sim 0.1\%$
\item Will conclusively test the coincidence
\end{itemize}

\subsection{Alternative Research Paths}

Given the null result for $\alpha_s$, two alternative directions are proposed:

\textbf{Path A: $m_H = 4\varphi^3 e^2$}
\begin{itemize}
\item Higgs mass: $m_H = 125.20 \pm 0.08$ GeV (PDG 2024)
\item Formula: $4\varphi^3 e^2 = 125.20$ GeV ($\Delta = 0.002\%$)
\item Question: Can this be derived algebraically?
\item Status: Sub-neighborhood density = 1 (unique at this precision)
\end{itemize}

\textbf{Path B: $m_e$ non-expressibility proof}
\begin{itemize}
\item Electron mass: $m_e = 0.51100$ MeV
\item Trinity basis: $\{\varphi, \pi, e\}$
\item Goal: Prove no $\{\varphi, \pi, e\}$ expression matches $m_e$ to $10^{-4}$ relative precision
\item Value: Negative result constrains Trinity framework
\end{itemize}

\subsection{Final Assessment}

The golden ratio expression for $\alpha_s$ represents:
\begin{itemize}
\item A \textbf{curious numerical coincidence} with $0.04\sigma$ agreement
\item An \textbf{algebraically elegant} result (7 steps from $\varphi^2 = \varphi + 1$)
\item A \textbf{mechanistically mysterious} relationship (no known theory)
\end{itemize}

\textbf{Position:} This paper does not claim a physical mechanism. Rather, it documents the coincidence and reports the failure to find a theoretical explanation. The expression stands as an open problem for theoretical physics.

\appendix

\section{50-Digit Seal}

\[
\alpha_\varphi = \varphi^{-3}/2 = (\sqrt{5} - 2)/2
\]

Numerical value (50 digits):
\begin{verbatim}
0.11803398874989484820458683436563811772030917980576
\end{verbatim}

Sealed for verification purposes. Any implementation claiming to test this formula must match this value to machine precision.

\textbf{Note:} This value was computed using Python's \texttt{mpmath} library with \texttt{prec=55} (55 decimal digits of precision). Standard Python float64 provides only 15-16 significant digits.

\section{Computational Verification}

All numerical results in Section IV were verified using:
\begin{itemize}
\item $E_8$ root system: exact computation in Python with mpmath (50-digit precision)
\item $H_4$ coordinates: verified against Coxeter group literature
\item Lie branching: checked with Lie algebra software (SAGE, LiE)
\end{itemize}

\textbf{Code repository:} \url{https://github.com/gHashTag/t27}

\textbf{Verification scripts:} \texttt{scripts/physics/verify\_e8\_phi.py}

\medskip
\noindent
\textbf{Status:} Pre-publication draft \\
\textbf{License:} MIT \\
\textbf{Correspondence:} t27 repository issues

\bibliographystyle{plain}
\bibliography{references}

\end{document}
